% from ega0-16-fr.tex
\oldpage[0\textsubscript{IV-1}]{23}
Si $S$ est une partie multiplicative de $A$, les idéaux premiers de $S^{-1}A$
correspondent biunivoquement aux idéaux premiers de $A$ ne rencontrant pas
$S$, donc on a
\[
\label{0.16.1.3.1-fr}
  \dim(S^{-1}A)\leq\dim(A).
\tag{16.1.3.1}
\]
Pour tout idéal premier $\mathfrak{p}$ de $A$, on a, par définition (14.2.1)
\[
\label{0.16.1.3.2-fr}
  \dim(A_{\mathfrak{p}})
  =\operatorname{codim}(V(\mathfrak{p}),\operatorname{Spec}(A)).
\tag{16.1.3.2}
\]
Ce nombre est encore appelé hauteur de l'idéal premier $\mathfrak{p}$ et noté
$\operatorname{ht}(\mathfrak{p})$. Plus généralement, pour tout idéal
$\mathfrak{a}$ de $A$, on appelle hauteur de $\mathfrak{a}$ le nombre
\[
\label{0.16.1.3.3-fr}
  \operatorname{ht}(\mathfrak{a})
  =\operatorname{codim}(V(\mathfrak{a}),\operatorname{Spec}(A)).
\tag{16.1.3.3}
\]
Si $(\mathfrak{p}_\lambda)$ est la famille des idéaux premiers minimaux parmi
ceux qui contiennent $\mathfrak{a}$, on a donc (14.2.1)
\[
\label{0.16.1.3.4-fr}
  \operatorname{ht}(\mathfrak{a})
  =\inf_\lambda(\operatorname{ht}(\mathfrak{p}_\lambda)).
\tag{16.1.3.4}
\]
On a évidemment
\[
\label{0.16.1.3.5-fr}
  \dim(A)=\sup_{\mathfrak{m}}(\dim(A_{\mathfrak{m}}))
\tag{16.1.3.5}
\]
où $\mathfrak{m}$ parcourt l'ensemble des idéaux maximaux de $A$.
\end{env}

\begin{env}[16.1.4]
\label{0.16.1.4-fr}
Pour tout idéal premier $\mathfrak{p}$ de $A$, on a (14.2.2.2)
\[
\label{0.16.1.4.1-fr}
  \dim(A_{\mathfrak{p}})+\dim(A/\mathfrak{p})\leq\dim A.
\tag{16.1.4.1}
\]
On dit que $A$ est caténaire (resp. équidimensionnel, équicodimensionnel,
biéquidimensionnel) lorsque $\operatorname{Spec}(A)$ est caténaire (resp.
équidimensionnel,
équicodimensionnel, biéquidimensionnel). Lorsque $A$ est noethérien et
biéquidimensionnel, les deux membres de (16.1.4.1) sont égaux pour tout idéal
premier de $A$ (14.3.5).

Pour tout idéal $\mathfrak{a}$ de $A$, les idéaux premiers de
$A/\mathfrak{a}$ correspondent biunivoquement aux idéaux premiers de $A$
contenant $\mathfrak{a}$, donc, si $A$ est caténaire, il en est de même de
$A/\mathfrak{a}$. De même, pour toute partie multiplicative $S$ de $A$, les
idéaux premiers de $S^{-1}A$ correspondent biunivoquement aux idéaux premiers
de $A$ ne rencontrant pas $S$, donc si $A$ est caténaire, il en est de même de
$S^{-1}A$.

En outre, comme les idéaux premiers de $A$ contenus dans un idéal premier
$\mathfrak{p}$ correspondent biunivoquement aux idéaux premiers de
$A_{\mathfrak{p}}$, pour que $A$ soit caténaire, il faut et il suffit que
$A_{\mathfrak{p}}$ le soit pour tout $\mathfrak{p}$. Dire que $A$ est
caténaire signifie donc ((14.3.2) et (16.1.3.2)) que pour tout couple d'idéaux
premiers $\mathfrak{p}$, $\mathfrak{q}$ de $A$ tels que
$\mathfrak{q}\subset\mathfrak{p}$, on a
\[
\label{0.16.1.4.2-fr}
  \dim(A_{\mathfrak{q}})
  +\dim(A_{\mathfrak{p}}/\mathfrak{q}A_{\mathfrak{p}})
  =\dim(A_{\mathfrak{p}}).
\tag{16.1.4.2}
\]
Si $A$ est biéquidimensionnel, il en est de même de $A_{\mathfrak{p}}$ pour
tout idéal premier $\mathfrak{p}$ de $A$.
\end{env}

\begin{proposition}[16.1.5]
\label{0.16.1.5-fr}
Soit $\rho:A\to B$ un homomorphisme d'anneaux faisant de $B$ une $A$-algèbre
entière. Alors on a $\dim(B)\leq\dim(A)$~; si de plus $\rho$ est injectif, on a
$\dim(B)=\dim(A)$.
\end{proposition}

\begin{proof}
